% Results subsection: Embedding Diagnostics
\subsection{Embedding Diagnostics}\label{sec:res_diagnostics}

Three statistical tests probed whether the 1920-dimensional embeddings
carry label-relevant signal (Table~\ref{tab:embedding-diagnostics}).
Discriminative information is limited to two of six datasets, and even
there it lacks the geometric organization that nearest-neighbor
classifiers require.

\paragraph{Mann--Whitney $U$.}
Only Desmoid (395/1920 dims, 21\%) and GIST (679/1920, 35\%) show
significant per-dimension separation after Benjamini--Hochberg
correction. CRLM, Lipo, Liver, and Melanoma yield zero significant
dimensions.

\paragraph{HSIC/CKA.}
Desmoid ($\text{CKA}{=}0.052$, $p{=}0.0012$) and GIST
($\text{CKA}{=}0.080$, $p{=}0.0001$) confirm weak but real global
dependence between features and labels; both CKA values fall in the
weak band ($0.02$--$0.1$). CRLM approaches significance ($p{=}0.083$)
without crossing it; the remaining three datasets show none ($p{>}0.15$).

\paragraph{$k$NN agreement.}
No dataset achieves strong local label coherence. Desmoid and GIST have
negative adjusted agreement at $k{=}5$ ($\kappa{=}{-}0.037$ and
${-}0.008$), meaning their signal is not geometrically localized. Lipo
is the sole exception ($\kappa{=}0.113$, $p{=}0.021$): a weak local
pattern that univariate tests and global kernels miss.

These findings complement the attention analysis
(Section~\ref{sec:res_attention}): the encoder \emph{attends} to
lesion regions but does not reliably \emph{encode} class-discriminative
features within them for most datasets. Where signal exists, it is
distributed across dimensions rather than organized into locally
coherent neighborhoods, limiting nearest-neighbor methods such as
LoCalPFN.
